\begin{helper}[Reduced configuration assembly]
Let $C$ be a Scheide structural configuration, let $D$ be a
\noderef{ScheideStructuralSelectionData}, and let $T'$ carry the associated
\noderef{ScheideStructuralReorderingData}.  Then $T'$ admits a reduced
Scheide structural configuration $R$ with $V(T')=V(T)$: every old base
vertex is in the new base, and every old block is either retained with its
color, both vertex identities, and both edge identities unchanged, or has
both vertices absorbed into the new base.
\end{helper}
\begin{proof}
Write $a=D.a$, $r=D.r$, $\upsilon=D.\upsilon$, and $\rho=D.\rho$.  Define
$p'=C.p+2a$ and $s'=r$.  The new base consists of the old base followed by
the two vertices of every $\upsilon$-block.  Its base edges are the
corresponding old base edges followed by the two exact edge identities of
every such block.  For $t\in\operatorname{Fin}(r)$, define the $t$-th new
block by
\[
 \gamma'(t)=C.\gamma(\rho(t)),\quad
 u'_j(t)=C.\operatorname{blockVertex}(\rho(t),j),\quad
 f'_j(t)=C.\operatorname{blockEdge}(\rho(t),j).
\]
The exact position clauses in
\noderef{ScheideStructuralReorderingData} give the four required base and
block position equations.  Its length clause, the equality $C.s=a+r$ in
\noderef{ScheideStructuralSelectionData}, and the old length equation give
$T'.\operatorname{length}+1=p'+2s'$.

The new block-color map is injective because both $C.\gamma$ and $\rho$
are injective.  Block attachment, block color, and a base missing witness
are inherited from the corresponding old $\rho(t)$-block; every old
missing witness lies in the unchanged first $C.p$ base positions.  For the
new base-color exclusion, an old nonfirst base edge is covered by the old
exclusion field of \noderef{ScheideStructuralConfiguration}.  A newly
absorbed base edge has color $C.\gamma(\upsilon(h))$.  Equality with a
retained color would, by injectivity of $C.\gamma$, identify
$\upsilon(h)$ with $\rho(t)$, contradicting the defining complement of
$\upsilon$.  Thus the exclusion holds on the enlarged base.

For exit coverage, take a new exit $v$.  Exit-set equality in
\noderef{ScheideStructuralReorderingData} makes it an old exit.  The
bijection $D.b$ gives a unique $t$ with underlying vertex $v$, and the
selected-missing field says that $\gamma'(t)$ is missing there.  The
no-isolated-vertex field is copied from $C$.

To reconstruct global maximality, unfold the old
\noderef{GlobalMaximumTriple}.  Its \noderef{KTriple} is unchanged, the
reordering datum supplies the new Tashkinov-tree assertion, and its
vertex-set equality transfers the cardinality equality with the Tashkinov
order.  Hence the preceding data form a new
\noderef{ScheideStructuralConfiguration} on $T'$.

We verify the three extra fields of
\noderef{ScheideReducedConfiguration}.  Exit-set equality and
$r=|\operatorname{ScheideExitVertices}(T,\varphi)|$ give the required
cardinality.  For a new exit $v$, use the inverse of $D.b$ to obtain $t$;
the selected-missing field supplies existence.  If another indexed new
color is missing at $v$, the selected-uniqueness field of
\noderef{ScheideStructuralSelectionData} makes its index equal to $t$.
Conversely, for every $t$, the exit underlying $D.b(t)$ is in the new exit
set and misses $\gamma'(t)$.  These are respectively exact uniqueness and
surjectivity onto the new block colors.

Call the resulting reduced configuration $R$.  The required vertex-set
equality is already a field of the reordering datum.  Each old base vertex
is the new base vertex at the same initial position.  Finally fix an old
block index $i$ and apply the unique partition in
\noderef{ScheideStructuralSelectionData}.  If $i=\rho(t)$, the definitions
of the retained block give equality of its color, its two vertex
identities, and its two exact edge identities with block $t$ of $R$.  If
$i=\upsilon(h)$, the two old vertices occur at the new base positions
$C.p+2h+j$.  Hence both lie in the image of $R.y$.  This proves precisely
the retained-block versus absorbed-block alternative without identifying
parallel edge identities.
\end{proof}
